% ============================================================================
%  gsd-guides-preamble.tex
%  Shared LaTeX preamble for the GSD User Guides Suite
%  Compile with: xelatex (3-pass), NOT pdflatex
%  License: CC BY-SA 4.0 — see \cclicense macro
%
%  Load order note: this file loads xcolor + tcolorbox BEFORE sourcing
%  gsd-color-scheme.tex, and then loads hyperref AFTER color definitions
%  are live. Documents only need to % ============================================================================
%  gsd-guides-preamble.tex
%  Shared LaTeX preamble for the GSD User Guides Suite
%  Compile with: xelatex (3-pass), NOT pdflatex
%  License: CC BY-SA 4.0 — see \cclicense macro
%
%  Load order note: this file loads xcolor + tcolorbox BEFORE sourcing
%  gsd-color-scheme.tex, and then loads hyperref AFTER color definitions
%  are live. Documents only need to % ============================================================================
%  gsd-guides-preamble.tex
%  Shared LaTeX preamble for the GSD User Guides Suite
%  Compile with: xelatex (3-pass), NOT pdflatex
%  License: CC BY-SA 4.0 — see \cclicense macro
%
%  Load order note: this file loads xcolor + tcolorbox BEFORE sourcing
%  gsd-color-scheme.tex, and then loads hyperref AFTER color definitions
%  are live. Documents only need to % ============================================================================
%  gsd-guides-preamble.tex
%  Shared LaTeX preamble for the GSD User Guides Suite
%  Compile with: xelatex (3-pass), NOT pdflatex
%  License: CC BY-SA 4.0 — see \cclicense macro
%
%  Load order note: this file loads xcolor + tcolorbox BEFORE sourcing
%  gsd-color-scheme.tex, and then loads hyperref AFTER color definitions
%  are live. Documents only need to \input{gsd-guides-preamble} in their
%  preamble (this file pulls in the color scheme transitively).
% ============================================================================

% ---------- Fonts ----------
\usepackage{fontspec}
\setmainfont{DejaVu Serif}[
  Ligatures=TeX,
  Scale=0.95,
  BoldFont={DejaVu Serif Bold},
  ItalicFont={DejaVu Serif Italic},
  BoldItalicFont={DejaVu Serif Bold Italic}
]
\newfontfamily\headingfont{DejaVu Sans}[Scale=0.95]
\newfontfamily\monofont{DejaVu Sans Mono}[Scale=0.88]
\setmonofont{DejaVu Sans Mono}[Scale=0.88]

% ---------- Page geometry ----------
\usepackage[
  letterpaper,
  top=1in,
  bottom=1in,
  left=1.25in,
  right=1.25in,
  headheight=14pt
]{geometry}

% ---------- Color + callout engine (must load BEFORE color-scheme input) ----
\usepackage{xcolor}
\usepackage[most]{tcolorbox}

% ---------- Typesetting packages ----------
\usepackage{tabularx}
\usepackage{longtable}
\usepackage{booktabs}
\usepackage{colortbl}
\usepackage{lastpage}
\usepackage{enumitem}
\usepackage{listings}
\usepackage{tikz}
\usepackage{fancyhdr}
\usepackage{titlesec}
\usepackage{parskip}
\usepackage{microtype}

% ---------- Pull in the color scheme (defines gsdnavy, gsdaccent, callouts) -
\input{gsd-color-scheme.tex}

% ---------- Hyperref (needs colors already defined) ----------
\usepackage[
  colorlinks=true,
  linkcolor=gsdnavy,
  urlcolor=gsdaccent,
  citecolor=gsdnavy,
  bookmarksnumbered=true
]{hyperref}

% ---------- Header / footer ----------
\pagestyle{fancy}
\fancyhf{}
\fancyhead[L]{\small\headingfont\@title}
\fancyhead[R]{\small\headingfont\leftmark}
\fancyfoot[L]{\small\headingfont GSD v1.35.0 — CC BY-SA 4.0}
\fancyfoot[R]{\small\headingfont Page \thepage\ of \pageref*{LastPage}}
\renewcommand{\headrulewidth}{0.4pt}
\renewcommand{\footrulewidth}{0.4pt}

\fancypagestyle{plain}{%
  \fancyhf{}
  \fancyfoot[L]{\small\headingfont GSD v1.35.0 — CC BY-SA 4.0}
  \fancyfoot[R]{\small\headingfont Page \thepage\ of \pageref*{LastPage}}
  \renewcommand{\headrulewidth}{0pt}
  \renewcommand{\footrulewidth}{0.4pt}
}

% ---------- Chapter / section styling ----------
\titleformat{\chapter}[display]
  {\headingfont\huge\bfseries\color{gsdnavy}}
  {\chaptertitlename\ \thechapter}{20pt}{\Huge}
\titleformat{\section}
  {\headingfont\Large\bfseries\color{gsdnavy}}
  {\thesection}{1em}{}
\titleformat{\subsection}
  {\headingfont\large\bfseries\color{gsdnavy}}
  {\thesubsection}{1em}{}

% ---------- Listings (code blocks) ----------
\lstdefinestyle{gsdcode}{
  basicstyle=\ttfamily\small,
  backgroundcolor=\color{gsdcodebg},
  frame=leftline,
  framerule=2pt,
  rulecolor=\color{gsdaccent},
  xleftmargin=1em,
  xrightmargin=0.5em,
  aboveskip=0.8em,
  belowskip=0.8em,
  breaklines=true,
  breakatwhitespace=true,
  showstringspaces=false,
  columns=flexible,
  keepspaces=true
}
\lstset{style=gsdcode}

% ---------- Semantic macros ----------
\newcommand{\cclicense}{%
  \textit{Released under CC BY-SA 4.0.\\
  Authored by Tibsfox. GSD itself is MIT-licensed — see gsd-build/get-shit-done.}%
}

\newcommand{\chapterheading}[1]{{\headingfont\LARGE\bfseries\color{gsdnavy}#1\par\medskip}}
\newcommand{\sectionheading}[1]{{\headingfont\Large\bfseries\color{gsdnavy}#1\par\smallskip}}
 in their
%  preamble (this file pulls in the color scheme transitively).
% ============================================================================

% ---------- Fonts ----------
\usepackage{fontspec}
\setmainfont{DejaVu Serif}[
  Ligatures=TeX,
  Scale=0.95,
  BoldFont={DejaVu Serif Bold},
  ItalicFont={DejaVu Serif Italic},
  BoldItalicFont={DejaVu Serif Bold Italic}
]
\newfontfamily\headingfont{DejaVu Sans}[Scale=0.95]
\newfontfamily\monofont{DejaVu Sans Mono}[Scale=0.88]
\setmonofont{DejaVu Sans Mono}[Scale=0.88]

% ---------- Page geometry ----------
\usepackage[
  letterpaper,
  top=1in,
  bottom=1in,
  left=1.25in,
  right=1.25in,
  headheight=14pt
]{geometry}

% ---------- Color + callout engine (must load BEFORE color-scheme input) ----
\usepackage{xcolor}
\usepackage[most]{tcolorbox}

% ---------- Typesetting packages ----------
\usepackage{tabularx}
\usepackage{longtable}
\usepackage{booktabs}
\usepackage{colortbl}
\usepackage{lastpage}
\usepackage{enumitem}
\usepackage{listings}
\usepackage{tikz}
\usepackage{fancyhdr}
\usepackage{titlesec}
\usepackage{parskip}
\usepackage{microtype}

% ---------- Pull in the color scheme (defines gsdnavy, gsdaccent, callouts) -
% ============================================================================
%  gsd-color-scheme.tex
%  Color palette + tcolorbox environment definitions for all GSD User Guides
% ============================================================================

\definecolor{gsdnavy}{HTML}{1B2A4A}
\definecolor{gsdaccent}{HTML}{CB3837}
\definecolor{gsdcodebg}{HTML}{F5F5F5}
\definecolor{gsdcodeborder}{HTML}{CCCCCC}

\definecolor{gsdtipbg}{HTML}{E8F4FD}
\definecolor{gsdtipbar}{HTML}{2B6CB0}

\definecolor{gsdwarnbg}{HTML}{FFF3CD}
\definecolor{gsdwarnbar}{HTML}{B78D00}

\definecolor{gsddangerbg}{HTML}{F8D7DA}
\definecolor{gsddangerbar}{HTML}{B02A37}

\definecolor{gsdsuccessbg}{HTML}{D4EDDA}
\definecolor{gsdsuccessbar}{HTML}{198754}

\definecolor{gsdnotebg}{HTML}{F0F0F5}
\definecolor{gsdnotebar}{HTML}{495057}

% ---------- Callout boxes ----------
\tcbset{
  gsdcalloutbase/.style={
    enhanced,
    breakable,
    left=10pt,
    right=8pt,
    top=6pt,
    bottom=6pt,
    boxrule=0pt,
    frame hidden,
    borderline west={4pt}{0pt}{#1},
    colback=#1,
    fonttitle=\bfseries\headingfont,
    fontupper=\small,
    arc=2pt
  }
}

\newtcolorbox{gsdtip}[1][Tip]{%
  gsdcalloutbase=gsdtipbg,
  borderline west={4pt}{0pt}{gsdtipbar},
  colback=gsdtipbg,
  title={#1}
}

\newtcolorbox{gsdwarning}[1][Warning]{%
  gsdcalloutbase=gsdwarnbg,
  borderline west={4pt}{0pt}{gsdwarnbar},
  colback=gsdwarnbg,
  title={#1}
}

\newtcolorbox{gsdimportant}[1][Important]{%
  gsdcalloutbase=gsddangerbg,
  borderline west={4pt}{0pt}{gsddangerbar},
  colback=gsddangerbg,
  title={#1}
}

\newtcolorbox{gsdnote}[1][Note]{%
  gsdcalloutbase=gsdnotebg,
  borderline west={4pt}{0pt}{gsdnotebar},
  colback=gsdnotebg,
  title={#1}
}

\newtcolorbox{gsdsuccess}[1][Success]{%
  gsdcalloutbase=gsdsuccessbg,
  borderline west={4pt}{0pt}{gsdsuccessbar},
  colback=gsdsuccessbg,
  title={#1}
}


% ---------- Hyperref (needs colors already defined) ----------
\usepackage[
  colorlinks=true,
  linkcolor=gsdnavy,
  urlcolor=gsdaccent,
  citecolor=gsdnavy,
  bookmarksnumbered=true
]{hyperref}

% ---------- Header / footer ----------
\pagestyle{fancy}
\fancyhf{}
\fancyhead[L]{\small\headingfont\@title}
\fancyhead[R]{\small\headingfont\leftmark}
\fancyfoot[L]{\small\headingfont GSD v1.35.0 — CC BY-SA 4.0}
\fancyfoot[R]{\small\headingfont Page \thepage\ of \pageref*{LastPage}}
\renewcommand{\headrulewidth}{0.4pt}
\renewcommand{\footrulewidth}{0.4pt}

\fancypagestyle{plain}{%
  \fancyhf{}
  \fancyfoot[L]{\small\headingfont GSD v1.35.0 — CC BY-SA 4.0}
  \fancyfoot[R]{\small\headingfont Page \thepage\ of \pageref*{LastPage}}
  \renewcommand{\headrulewidth}{0pt}
  \renewcommand{\footrulewidth}{0.4pt}
}

% ---------- Chapter / section styling ----------
\titleformat{\chapter}[display]
  {\headingfont\huge\bfseries\color{gsdnavy}}
  {\chaptertitlename\ \thechapter}{20pt}{\Huge}
\titleformat{\section}
  {\headingfont\Large\bfseries\color{gsdnavy}}
  {\thesection}{1em}{}
\titleformat{\subsection}
  {\headingfont\large\bfseries\color{gsdnavy}}
  {\thesubsection}{1em}{}

% ---------- Listings (code blocks) ----------
\lstdefinestyle{gsdcode}{
  basicstyle=\ttfamily\small,
  backgroundcolor=\color{gsdcodebg},
  frame=leftline,
  framerule=2pt,
  rulecolor=\color{gsdaccent},
  xleftmargin=1em,
  xrightmargin=0.5em,
  aboveskip=0.8em,
  belowskip=0.8em,
  breaklines=true,
  breakatwhitespace=true,
  showstringspaces=false,
  columns=flexible,
  keepspaces=true
}
\lstset{style=gsdcode}

% ---------- Semantic macros ----------
\newcommand{\cclicense}{%
  \textit{Released under CC BY-SA 4.0.\\
  Authored by Tibsfox. GSD itself is MIT-licensed — see gsd-build/get-shit-done.}%
}

\newcommand{\chapterheading}[1]{{\headingfont\LARGE\bfseries\color{gsdnavy}#1\par\medskip}}
\newcommand{\sectionheading}[1]{{\headingfont\Large\bfseries\color{gsdnavy}#1\par\smallskip}}
 in their
%  preamble (this file pulls in the color scheme transitively).
% ============================================================================

% ---------- Fonts ----------
\usepackage{fontspec}
\setmainfont{DejaVu Serif}[
  Ligatures=TeX,
  Scale=0.95,
  BoldFont={DejaVu Serif Bold},
  ItalicFont={DejaVu Serif Italic},
  BoldItalicFont={DejaVu Serif Bold Italic}
]
\newfontfamily\headingfont{DejaVu Sans}[Scale=0.95]
\newfontfamily\monofont{DejaVu Sans Mono}[Scale=0.88]
\setmonofont{DejaVu Sans Mono}[Scale=0.88]

% ---------- Page geometry ----------
\usepackage[
  letterpaper,
  top=1in,
  bottom=1in,
  left=1.25in,
  right=1.25in,
  headheight=14pt
]{geometry}

% ---------- Color + callout engine (must load BEFORE color-scheme input) ----
\usepackage{xcolor}
\usepackage[most]{tcolorbox}

% ---------- Typesetting packages ----------
\usepackage{tabularx}
\usepackage{longtable}
\usepackage{booktabs}
\usepackage{colortbl}
\usepackage{lastpage}
\usepackage{enumitem}
\usepackage{listings}
\usepackage{tikz}
\usepackage{fancyhdr}
\usepackage{titlesec}
\usepackage{parskip}
\usepackage{microtype}

% ---------- Pull in the color scheme (defines gsdnavy, gsdaccent, callouts) -
% ============================================================================
%  gsd-color-scheme.tex
%  Color palette + tcolorbox environment definitions for all GSD User Guides
% ============================================================================

\definecolor{gsdnavy}{HTML}{1B2A4A}
\definecolor{gsdaccent}{HTML}{CB3837}
\definecolor{gsdcodebg}{HTML}{F5F5F5}
\definecolor{gsdcodeborder}{HTML}{CCCCCC}

\definecolor{gsdtipbg}{HTML}{E8F4FD}
\definecolor{gsdtipbar}{HTML}{2B6CB0}

\definecolor{gsdwarnbg}{HTML}{FFF3CD}
\definecolor{gsdwarnbar}{HTML}{B78D00}

\definecolor{gsddangerbg}{HTML}{F8D7DA}
\definecolor{gsddangerbar}{HTML}{B02A37}

\definecolor{gsdsuccessbg}{HTML}{D4EDDA}
\definecolor{gsdsuccessbar}{HTML}{198754}

\definecolor{gsdnotebg}{HTML}{F0F0F5}
\definecolor{gsdnotebar}{HTML}{495057}

% ---------- Callout boxes ----------
\tcbset{
  gsdcalloutbase/.style={
    enhanced,
    breakable,
    left=10pt,
    right=8pt,
    top=6pt,
    bottom=6pt,
    boxrule=0pt,
    frame hidden,
    borderline west={4pt}{0pt}{#1},
    colback=#1,
    fonttitle=\bfseries\headingfont,
    fontupper=\small,
    arc=2pt
  }
}

\newtcolorbox{gsdtip}[1][Tip]{%
  gsdcalloutbase=gsdtipbg,
  borderline west={4pt}{0pt}{gsdtipbar},
  colback=gsdtipbg,
  title={#1}
}

\newtcolorbox{gsdwarning}[1][Warning]{%
  gsdcalloutbase=gsdwarnbg,
  borderline west={4pt}{0pt}{gsdwarnbar},
  colback=gsdwarnbg,
  title={#1}
}

\newtcolorbox{gsdimportant}[1][Important]{%
  gsdcalloutbase=gsddangerbg,
  borderline west={4pt}{0pt}{gsddangerbar},
  colback=gsddangerbg,
  title={#1}
}

\newtcolorbox{gsdnote}[1][Note]{%
  gsdcalloutbase=gsdnotebg,
  borderline west={4pt}{0pt}{gsdnotebar},
  colback=gsdnotebg,
  title={#1}
}

\newtcolorbox{gsdsuccess}[1][Success]{%
  gsdcalloutbase=gsdsuccessbg,
  borderline west={4pt}{0pt}{gsdsuccessbar},
  colback=gsdsuccessbg,
  title={#1}
}


% ---------- Hyperref (needs colors already defined) ----------
\usepackage[
  colorlinks=true,
  linkcolor=gsdnavy,
  urlcolor=gsdaccent,
  citecolor=gsdnavy,
  bookmarksnumbered=true
]{hyperref}

% ---------- Header / footer ----------
\pagestyle{fancy}
\fancyhf{}
\fancyhead[L]{\small\headingfont\@title}
\fancyhead[R]{\small\headingfont\leftmark}
\fancyfoot[L]{\small\headingfont GSD v1.35.0 — CC BY-SA 4.0}
\fancyfoot[R]{\small\headingfont Page \thepage\ of \pageref*{LastPage}}
\renewcommand{\headrulewidth}{0.4pt}
\renewcommand{\footrulewidth}{0.4pt}

\fancypagestyle{plain}{%
  \fancyhf{}
  \fancyfoot[L]{\small\headingfont GSD v1.35.0 — CC BY-SA 4.0}
  \fancyfoot[R]{\small\headingfont Page \thepage\ of \pageref*{LastPage}}
  \renewcommand{\headrulewidth}{0pt}
  \renewcommand{\footrulewidth}{0.4pt}
}

% ---------- Chapter / section styling ----------
\titleformat{\chapter}[display]
  {\headingfont\huge\bfseries\color{gsdnavy}}
  {\chaptertitlename\ \thechapter}{20pt}{\Huge}
\titleformat{\section}
  {\headingfont\Large\bfseries\color{gsdnavy}}
  {\thesection}{1em}{}
\titleformat{\subsection}
  {\headingfont\large\bfseries\color{gsdnavy}}
  {\thesubsection}{1em}{}

% ---------- Listings (code blocks) ----------
\lstdefinestyle{gsdcode}{
  basicstyle=\ttfamily\small,
  backgroundcolor=\color{gsdcodebg},
  frame=leftline,
  framerule=2pt,
  rulecolor=\color{gsdaccent},
  xleftmargin=1em,
  xrightmargin=0.5em,
  aboveskip=0.8em,
  belowskip=0.8em,
  breaklines=true,
  breakatwhitespace=true,
  showstringspaces=false,
  columns=flexible,
  keepspaces=true
}
\lstset{style=gsdcode}

% ---------- Semantic macros ----------
\newcommand{\cclicense}{%
  \textit{Released under CC BY-SA 4.0.\\
  Authored by Tibsfox. GSD itself is MIT-licensed — see gsd-build/get-shit-done.}%
}

\newcommand{\chapterheading}[1]{{\headingfont\LARGE\bfseries\color{gsdnavy}#1\par\medskip}}
\newcommand{\sectionheading}[1]{{\headingfont\Large\bfseries\color{gsdnavy}#1\par\smallskip}}
 in their
%  preamble (this file pulls in the color scheme transitively).
% ============================================================================

% ---------- Fonts ----------
\usepackage{fontspec}
\setmainfont{DejaVu Serif}[
  Ligatures=TeX,
  Scale=0.95,
  BoldFont={DejaVu Serif Bold},
  ItalicFont={DejaVu Serif Italic},
  BoldItalicFont={DejaVu Serif Bold Italic}
]
\newfontfamily\headingfont{DejaVu Sans}[Scale=0.95]
\newfontfamily\monofont{DejaVu Sans Mono}[Scale=0.88]
\setmonofont{DejaVu Sans Mono}[Scale=0.88]

% ---------- Page geometry ----------
\usepackage[
  letterpaper,
  top=1in,
  bottom=1in,
  left=1.25in,
  right=1.25in,
  headheight=14pt
]{geometry}

% ---------- Color + callout engine (must load BEFORE color-scheme input) ----
\usepackage{xcolor}
\usepackage[most]{tcolorbox}

% ---------- Typesetting packages ----------
\usepackage{tabularx}
\usepackage{longtable}
\usepackage{booktabs}
\usepackage{colortbl}
\usepackage{lastpage}
\usepackage{enumitem}
\usepackage{listings}
\usepackage{tikz}
\usepackage{fancyhdr}
\usepackage{titlesec}
\usepackage{parskip}
\usepackage{microtype}

% ---------- Pull in the color scheme (defines gsdnavy, gsdaccent, callouts) -
% ============================================================================
%  gsd-color-scheme.tex
%  Color palette + tcolorbox environment definitions for all GSD User Guides
% ============================================================================

\definecolor{gsdnavy}{HTML}{1B2A4A}
\definecolor{gsdaccent}{HTML}{CB3837}
\definecolor{gsdcodebg}{HTML}{F5F5F5}
\definecolor{gsdcodeborder}{HTML}{CCCCCC}

\definecolor{gsdtipbg}{HTML}{E8F4FD}
\definecolor{gsdtipbar}{HTML}{2B6CB0}

\definecolor{gsdwarnbg}{HTML}{FFF3CD}
\definecolor{gsdwarnbar}{HTML}{B78D00}

\definecolor{gsddangerbg}{HTML}{F8D7DA}
\definecolor{gsddangerbar}{HTML}{B02A37}

\definecolor{gsdsuccessbg}{HTML}{D4EDDA}
\definecolor{gsdsuccessbar}{HTML}{198754}

\definecolor{gsdnotebg}{HTML}{F0F0F5}
\definecolor{gsdnotebar}{HTML}{495057}

% ---------- Callout boxes ----------
\tcbset{
  gsdcalloutbase/.style={
    enhanced,
    breakable,
    left=10pt,
    right=8pt,
    top=6pt,
    bottom=6pt,
    boxrule=0pt,
    frame hidden,
    borderline west={4pt}{0pt}{#1},
    colback=#1,
    fonttitle=\bfseries\headingfont,
    fontupper=\small,
    arc=2pt
  }
}

\newtcolorbox{gsdtip}[1][Tip]{%
  gsdcalloutbase=gsdtipbg,
  borderline west={4pt}{0pt}{gsdtipbar},
  colback=gsdtipbg,
  title={#1}
}

\newtcolorbox{gsdwarning}[1][Warning]{%
  gsdcalloutbase=gsdwarnbg,
  borderline west={4pt}{0pt}{gsdwarnbar},
  colback=gsdwarnbg,
  title={#1}
}

\newtcolorbox{gsdimportant}[1][Important]{%
  gsdcalloutbase=gsddangerbg,
  borderline west={4pt}{0pt}{gsddangerbar},
  colback=gsddangerbg,
  title={#1}
}

\newtcolorbox{gsdnote}[1][Note]{%
  gsdcalloutbase=gsdnotebg,
  borderline west={4pt}{0pt}{gsdnotebar},
  colback=gsdnotebg,
  title={#1}
}

\newtcolorbox{gsdsuccess}[1][Success]{%
  gsdcalloutbase=gsdsuccessbg,
  borderline west={4pt}{0pt}{gsdsuccessbar},
  colback=gsdsuccessbg,
  title={#1}
}


% ---------- Hyperref (needs colors already defined) ----------
\usepackage[
  colorlinks=true,
  linkcolor=gsdnavy,
  urlcolor=gsdaccent,
  citecolor=gsdnavy,
  bookmarksnumbered=true
]{hyperref}

% ---------- Header / footer ----------
\pagestyle{fancy}
\fancyhf{}
\fancyhead[L]{\small\headingfont\@title}
\fancyhead[R]{\small\headingfont\leftmark}
\fancyfoot[L]{\small\headingfont GSD v1.35.0 — CC BY-SA 4.0}
\fancyfoot[R]{\small\headingfont Page \thepage\ of \pageref*{LastPage}}
\renewcommand{\headrulewidth}{0.4pt}
\renewcommand{\footrulewidth}{0.4pt}

\fancypagestyle{plain}{%
  \fancyhf{}
  \fancyfoot[L]{\small\headingfont GSD v1.35.0 — CC BY-SA 4.0}
  \fancyfoot[R]{\small\headingfont Page \thepage\ of \pageref*{LastPage}}
  \renewcommand{\headrulewidth}{0pt}
  \renewcommand{\footrulewidth}{0.4pt}
}

% ---------- Chapter / section styling ----------
\titleformat{\chapter}[display]
  {\headingfont\huge\bfseries\color{gsdnavy}}
  {\chaptertitlename\ \thechapter}{20pt}{\Huge}
\titleformat{\section}
  {\headingfont\Large\bfseries\color{gsdnavy}}
  {\thesection}{1em}{}
\titleformat{\subsection}
  {\headingfont\large\bfseries\color{gsdnavy}}
  {\thesubsection}{1em}{}

% ---------- Listings (code blocks) ----------
\lstdefinestyle{gsdcode}{
  basicstyle=\ttfamily\small,
  backgroundcolor=\color{gsdcodebg},
  frame=leftline,
  framerule=2pt,
  rulecolor=\color{gsdaccent},
  xleftmargin=1em,
  xrightmargin=0.5em,
  aboveskip=0.8em,
  belowskip=0.8em,
  breaklines=true,
  breakatwhitespace=true,
  showstringspaces=false,
  columns=flexible,
  keepspaces=true
}
\lstset{style=gsdcode}

% ---------- Semantic macros ----------
\newcommand{\cclicense}{%
  \textit{Released under CC BY-SA 4.0.\\
  Authored by Tibsfox. GSD itself is MIT-licensed — see gsd-build/get-shit-done.}%
}

\newcommand{\chapterheading}[1]{{\headingfont\LARGE\bfseries\color{gsdnavy}#1\par\medskip}}
\newcommand{\sectionheading}[1]{{\headingfont\Large\bfseries\color{gsdnavy}#1\par\smallskip}}
